\documentclass[11pt,a4paper]{article}
\usepackage[utf8]{inputenc}
\usepackage[T1]{fontenc}
\usepackage{amsmath,amssymb,amsfonts}
\usepackage{booktabs}
\usepackage{hyperref}
\usepackage{graphicx}
\usepackage{algorithm}
\usepackage{algorithmic}
\usepackage[margin=1in]{geometry}
\usepackage{natbib}
\usepackage{xcolor}

\title{OpenProof-Tactic: On-Device Lean 4 Tactic Prediction\\via Hybrid-Architecture Small Language Models}

\author{
Mark Miller \\
OpenProof \\
\texttt{mark@openproof.ai}
}

\date{March 2026}

\begin{document}

\maketitle

\begin{abstract}
We present OpenProof-Tactic, a small on-device model for step-level Lean 4 tactic prediction. We compare two base models: \textbf{Qwen3.5-2B}, the newest hybrid Gated DeltaNet architecture (released March 2026), and \textbf{Qwen3-1.7B}, a proven standard transformer base (used by Kimina-Prover-RL-1.7B at 76.6\% pass@32 on MiniF2F). Both are trained in three stages: (1) SFT on 1.2M (state, tactic) pairs extracted via Pantograph from Mathlib4, Lean Workbook, and Goedel-Pset; (2) three rounds of expert iteration through self-play proof search; and (3) GRPO reinforcement learning with binary Lean compiler feedback. The Qwen3.5-2B variant achieves \textbf{XX.X\%} pass@32 on MiniF2F-test, surpassing the prior small-model SOTA (Kimina-Prover-RL-1.7B, 76.6\%) and approaching 7B-class systems. At Q4 quantization ($\sim$1.5GB), the model generates tactics in under 300ms on Apple Silicon with no network dependency. We release weights, training data, and code.
\end{abstract}

%----------------------------------------------------------------------
\section{Introduction}
\label{sec:intro}

Automated theorem proving in Lean 4 has progressed from 26\% on MiniF2F-test (ReProver, 2023) to 99.2\% (HILBERT, 2025) in under two years. But all competitive systems require either large models (7B+, 5--8GB RAM) or cloud inference. This creates a barrier for mathematicians who want fast, private, offline proof assistance.

Three developments make a small on-device tactic model newly tractable:

\begin{enumerate}
\item \textbf{Step-level search works.} BFS-Prover-V2 \citep{bfsprover2} showed that best-first search with a step-level model beats complex MCTS and whole-proof generation. The search handles proof structure; the model just suggests one tactic per goal state.

\item \textbf{Small models have proven sufficient.} Kimina-Prover-RL-1.7B \citep{kimina} achieves 76.6\% pass@32 on MiniF2F at just 1.7B parameters---within striking distance of 7B systems. The key is expert iteration and RL, not raw model scale.

\item \textbf{Base model quality is rapidly improving.} Qwen3.5 \citep{qwen35} (March 2026) introduces hybrid Gated DeltaNet attention, achieving 2B-parameter performance that matches prior 7B models on reasoning benchmarks. No one has yet applied these newest architectures to theorem proving.
\end{enumerate}

Our central hypothesis is that \textbf{the newest general-purpose small models, trained with the right data and RL pipeline, can match or exceed math-specialized models at tactic prediction}. We test this by comparing:

\begin{itemize}
\item \textbf{Qwen3.5-2B}: The newest architecture (March 2, 2026). Hybrid GDN + softmax attention, 262K context, 2B params. Never before used for theorem proving. Potentially the strongest base at this size, but unproven for the domain and has known quantization caveats.
\item \textbf{Qwen3-1.7B}: Standard transformer (April 2025). 40K context, 1.7B params. Proven base---Kimina-Prover-Distill-1.7B achieves 73\% and Kimina-Prover-RL-1.7B achieves 76.6\% on MiniF2F. The established safe choice.
\end{itemize}

Both are trained with the same three-stage pipeline (SFT, expert iteration, GRPO) on identical data, isolating the effect of the base model.

%----------------------------------------------------------------------
\section{Related Work}
\label{sec:related}

\paragraph{Step-level tactic models.} LeanDojo/ReProver \citep{leandojo} introduced retrieval-augmented tactic prediction (ByT5, $\sim$26\%). llmstep \citep{llmstep} fine-tuned Pythia-2.8B on LeanDojo data (27.9\%). BFS-Prover \citep{bfsprover1} scaled to 73\% with expert iteration and DPO on Qwen2.5-Math-7B. BFS-Prover-V2 \citep{bfsprover2} reached 82.4\% (7B) and 95\% (32B + planner) via multi-stage expert iteration. All use the same format: goal state in, tactic out.

\paragraph{Small prover models.} Kimina-Prover \citep{kimina} distilled a 72B RL-trained model to 1.7B, with the distilled variant achieving 73\% pass@32 and the RL-refined variant 76.6\% pass@32 on MiniF2F. This is the current small-model SOTA. Kimina-Prover-Preview-Distill-1.5B on the older Qwen2.5-Math-1.5B base reaches 42.6\% pass@1. The gap between the Qwen2.5-Math-1.5B variant (42.6\% p@1) and the Qwen3-1.7B variant (76.6\% p@32 with RL) demonstrates that both the base model and the training recipe matter.

\paragraph{Whole-proof generation.} DeepSeek-Prover-V2 \citep{deepseekprover2} (88.9\%), Goedel-Prover-V2 \citep{goedel} (90.4\% with self-correction), and Kimina-Prover-72B (80.7\%) generate complete proofs with chain-of-thought. Powerful but require 7B+ parameters and are less amenable to tree search.

\paragraph{Agentic systems.} HILBERT \citep{hilbert} (99.2\%), Seed-Prover 1.5 \citep{seedprover15}, and Delta-Prover \citep{deltaprover} (95.9\%) wrap frontier LLMs with decomposition, retrieval, and repair. Our tactic model improves the inner search loop of any such system.

\paragraph{Training techniques.} GRPO \citep{deepseekmath} provides critic-free RL; DAPO \citep{dapo} improves it with asymmetric clipping and dynamic sampling for binary-reward settings. V-STaR \citep{vstar} extends self-play by training on both successes and failures. TTRL \citep{ttrl} enables self-improvement on unlabeled data. Prompt Augmentation for GRPO \citep{promptauggrpo} stabilizes training for small models. RLAD \citep{rlad} improves RL-aware distillation. 
\paragraph{Hybrid architectures.} Qwen3.5 \citep{qwen35} uses Gated DeltaNet (GDN) linear attention for 75\% of layers and full softmax attention for 25\%, achieving O(1) per-token inference in the linear layers. Falcon-H1 \citep{falconh1} uses Mamba-2. Neither has been applied to formal theorem proving. We are the first to evaluate GDN-based models for tactic prediction.

%----------------------------------------------------------------------
\section{Method}
\label{sec:method}

\subsection{Base Models}
\label{sec:base}

We train two variants from the \textbf{base} (non-instruct) checkpoints to avoid chat-template overhead:

\begin{itemize}
\item \textbf{Qwen3.5-2B-Base} \citep{qwen35}: 2B parameters, hybrid GDN + softmax architecture (3:1 ratio), 262K native context, released March 2, 2026. On general benchmarks, Qwen3.5-2B matches Falcon-H1R-7B and exceeds Nemotron Nano 9B V2 \citep{qwen35}. Q4\_K\_M quantization yields $\sim$1.5GB. \textbf{Caveat}: Unsloth advises against QLoRA for Qwen3.5 due to quantization sensitivity of GDN layers; we use LoRA at bf16 for fine-tuning.

\item \textbf{Qwen3-1.7B-Base} \citep{qwen3}: 1.7B parameters, standard GQA transformer, 40K context, released April 28, 2025. Matches Qwen2.5-3B on STEM/coding despite being smaller. 36T pretraining tokens. Q4\_K\_M yields $\sim$1GB. Proven for theorem proving: Kimina-Prover-RL-1.7B uses this base and achieves 76.6\% pass@32 on MiniF2F.
\end{itemize}

\subsection{Input Format}

We follow BFS-Prover-V2's format: raw completion with prompt \texttt{\{tactic\_state\}:::} and the model generates a single tactic, terminated by newline or \texttt{:::}. This is minimal, fast to tokenize, and directly compatible with OpenProof's search pipeline. No chat template or system prompt.

\subsection{Stage 1: Data Extraction and SFT}
\label{sec:stage1}

\paragraph{Data sources.} We extract (state, tactic) pairs from three sources:

\begin{itemize}
\item \textbf{Mathlib4 via LeanDojo Benchmark 4} \citep{leandojo}: $\sim$200K tactic steps from 98K theorems, replayed through Pantograph \citep{pantograph} to extract exact goal states in \texttt{target.pp} format.

\item \textbf{Lean Workbook Plus} \citep{leanworkbook}: $\sim$100K tactic steps from 25K solved competition-level problems.

\item \textbf{Goedel-Pset-v1-solved} \citep{goedel}: $\sim$900K tactic steps from 800K+ proofs of 1.78M formalized statements.
\end{itemize}

Total: \textbf{1.2M (state, tactic) pairs}. Each pair records the Pantograph \texttt{target.pp} goal state before tactic application. We discard \texttt{sorry}, \texttt{admit}, and \texttt{native\_decide}.

\paragraph{SFT training.}
\begin{itemize}
\item \textbf{Qwen3.5-2B}: LoRA at bf16 (rank 64, alpha 128). Full fine-tuning risks destabilizing GDN layers; LoRA preserves pretrained representations while adapting to the domain.
\item \textbf{Qwen3-1.7B}: Full fine-tuning (1.7B fits in A100-80GB with room for gradient accumulation).
\item Common: LR $2 \times 10^{-5}$ cosine decay to $1 \times 10^{-6}$, batch size 128, 3 epochs, max seq 2048.
\item Hardware: 1$\times$ A100-80GB. Training time: $\sim$6--8 hours per variant.
\end{itemize}

\subsection{Stage 2: Expert Iteration with V-STaR}
\label{sec:stage2}

We extend standard expert iteration \citep{expertiteration} with V-STaR \citep{vstar}: training on both successful \emph{and} failed proof attempts (with labels), which produces a more calibrated model and better search behavior.

\begin{algorithm}
\caption{Expert Iteration with V-STaR and Proof Simplification}
\begin{algorithmic}[1]
\STATE $\mathcal{M}_0 \leftarrow$ SFT model from Stage 1
\STATE $\mathcal{D}_0^+ \leftarrow$ initial successful (state, tactic) pairs
\STATE $\mathcal{D}_0^- \leftarrow$ failed (state, tactic) pairs from SFT training errors
\FOR{$i = 1$ to $3$}
  \FOR{each unsolved problem $p$ in corpus}
    \STATE Run best-first search with $\mathcal{M}_{i-1}$, Pantograph, beam=16, max=400
    \IF{proof found}
      \STATE Simplify proof (remove redundant tactics, re-verify)
      \STATE Add (state, tactic) pairs to $\mathcal{D}_i^+$
    \ENDIF
    \STATE Collect failed tactic attempts into $\mathcal{D}_i^-$
  \ENDFOR
  \STATE $\mathcal{M}_i \leftarrow$ SFT($\mathcal{M}_0$, $\bigcup_{j \leq i} \mathcal{D}_j^+$)
  \STATE $\mathcal{M}_i \leftarrow$ DPO($\mathcal{M}_i$, $\mathcal{D}_i^+$, $\mathcal{D}_i^-$) \COMMENT{V-STaR: learn from failures}
  \STATE Remove problems solvable at beam=1 (adaptive filtering \citep{bfsprover2})
\ENDFOR
\end{algorithmic}
\end{algorithm}

\paragraph{Problem corpus.} Lean Workbook Plus unsolved ($\sim$58K), Goedel-Pset unsolved ($\sim$980K), NuminaMath autoformalized ($\sim$860K). Total: $\sim$1.9M candidate problems.

\paragraph{Search config.} Beam 16, max 400 expansions, 120s timeout, length-normalized scoring ($\alpha = 0.1$), Pantograph at $\sim$3ms/tactic. We use \textbf{adaptive search budgets}: try greedy first (1 expansion), escalate to beam=4 (32 expansions), then full beam=16 (400 expansions). This allocates compute where it matters and accelerates the search phase by 3--5$\times$.

\paragraph{Expected yield.} Based on BFS-Prover-V2 ($\sim$10$^7$ searches/round), each round adds 100K--400K new positive pairs plus 500K--1M negative pairs. The V-STaR DPO step uses both.

\subsection{Stage 3: DAPO Reinforcement Learning}
\label{sec:stage3}

We use DAPO \citep{dapo} rather than GRPO, incorporating three key improvements for binary-verifier settings:

\paragraph{Per-tactic rewards.} Rather than a binary whole-proof signal, Pantograph provides \emph{step-level} verification. We define a shaped reward:

\[
r(\text{state}, \text{tactic}) = \begin{cases}
+1.0 & \text{tactic closes a goal (goals decrease)} \\
+0.5 & \text{tactic succeeds and changes the goal state} \\
\phantom{+}0.0 & \text{tactic produces a Lean error}
\end{cases}
\]

This dense reward signal is especially important for small models, where binary rewards produce high-variance gradients.

\paragraph{DAPO improvements over GRPO.}
\begin{itemize}
\item \textbf{Asymmetric clipping}: $\epsilon_{\text{low}} = 0.1$, $\epsilon_{\text{high}} = 0.2$. Wider upside clipping encourages exploration of novel tactics.
\item \textbf{No KL penalty}: with a perfect verifier, we do not need to anchor to a reference policy---the verifier prevents degenerate outputs.
\item \textbf{Dynamic sampling}: skip prompts where all $k$ rollouts succeed (too easy) or all fail (too hard). This focuses gradient signal on the model's ``frontier'' problems.
\item \textbf{Dr.\ GRPO length normalization} \citep{drgrpo}: normalize per-token advantages by sequence length to correct the length bias in standard GRPO.
\end{itemize}

\paragraph{Training.}
\begin{itemize}
\item Prompts: $\sim$5K unique theorem goal states, selected via difficulty curriculum
\item Rollouts: 32 per prompt
\item \textbf{Qwen3.5-2B}: LoRA at bf16 (consistent with SFT stage)
\item \textbf{Qwen3-1.7B}: QLoRA (DAPO is memory-intensive; QLoRA enables 32 rollouts per batch)
\item Framework: TRL \citep{trl} with Prompt Augmentation \citep{promptauggrpo} to prevent entropy collapse.
\item Hardware: 1$\times$ A100-80GB. Time: $\sim$12--24 hours.
\end{itemize}

\subsection{Quantization and Deployment}
\label{sec:deploy}

We quantize to Q4\_K\_M via llama.cpp and serve via ollama:

\begin{itemize}
\item \textbf{Qwen3.5-2B}: $\sim$1.5GB at Q4\_K\_M. Special care for GDN layers---we use Q4\_K (not MXFP4) for \texttt{ssm\_out}, \texttt{attn\_gate}, \texttt{ssm\_beta}, \texttt{ssm\_alpha} tensors.
\item \textbf{Qwen3-1.7B}: $\sim$1GB at Q4\_K\_M. Standard GGUF quantization, maximally compatible.
\end{itemize}

Integration with OpenProof: zero search code changes. The existing \texttt{OllamaProposer} calls the model via HTTP; the \texttt{ProposeFn} interface passes goal states as strings.

%----------------------------------------------------------------------
\section{Experimental Setup}
\label{sec:experiments}

\subsection{Evaluation}

\paragraph{Benchmark.} MiniF2F-test \citep{minif2f}: 244 problems from AIME, AMC, IMO, and Mathlib, formalized in Lean 4.

\paragraph{Search.} Best-first with Pantograph: beam=8, max expansions=200, timeout=120s, dedup by goal hash, max depth=20, length penalty $\alpha{=}0.1$.

\paragraph{Metrics.} pass@1 (greedy), pass@32 (32 search attempts per goal), cumulative (full 200-expansion budget).

\subsection{Baselines}

\begin{table}[h]
\centering
\caption{Baseline models for comparison.}
\begin{tabular}{lccc}
\toprule
\textbf{Model} & \textbf{Params} & \textbf{MiniF2F} & \textbf{Type} \\
\midrule
Standard tactics & -- & 14.0\% cumul. & Rule-based \\
llmstep Pythia & 2.8B & 27.9\% cumul. & Step-level SFT \\
Kimina-Distill-1.5B & 1.5B & 42.6\% p@1 & Distillation \\
Kimina-Distill-1.7B & 1.7B & 73.0\% p@32 & Distillation \\
Kimina-RL-1.7B & 1.7B & 76.6\% p@32 & Distillation + RL \\
\midrule
BFS-Prover-V2-7B & 7B & 82.4\% cumul. & Expert iter. + RL \\
DeepSeek-Prover-V2-7B & 7B & 88.9\% & Subgoal decomp. \\
\bottomrule
\end{tabular}
\label{tab:baselines}
\end{table}

Our target: exceed Kimina-RL-1.7B (76.6\%) and narrow the gap to 7B systems.

\subsection{Ablations}

\begin{enumerate}
\item Base model comparison: Qwen3.5-2B vs Qwen3-1.7B (same data, same pipeline)
\item Training stage ablation: SFT only $\rightarrow$ + expert iteration $\rightarrow$ + GRPO $\rightarrow$ + DPO
\item Quantization: F16, Q8, Q4\_K\_M
\item Data scale: 200K (LeanDojo only) vs 1.2M (full)
\end{enumerate}

%----------------------------------------------------------------------
\section{Results}
\label{sec:results}

\subsection{Base Model Comparison}

\begin{table}[h]
\centering
\caption{Base model comparison at each training stage. MiniF2F-test pass@32.}
\begin{tabular}{lcc}
\toprule
\textbf{Training Stage} & \textbf{Qwen3.5-2B} & \textbf{Qwen3-1.7B} \\
\midrule
SFT (1.2M pairs) & XX.X\% & XX.X\% \\
+ Expert Iter. $\times$3 & XX.X\% & XX.X\% \\
+ GRPO & XX.X\% & XX.X\% \\
+ DPO & XX.X\% & XX.X\% \\
\bottomrule
\end{tabular}
\label{tab:base_model}
\end{table}

\subsection{Main Results}

\begin{table}[h]
\centering
\caption{MiniF2F-test results. All use beam=8, expansions=200 search. ``On-device'' = Q4, $\leq$2GB RAM.}
\begin{tabular}{lcccccc}
\toprule
\textbf{Model} & \textbf{Params} & \textbf{Base} & \textbf{p@1} & \textbf{p@32} & \textbf{Cumul.} & \textbf{On-dev.} \\
\midrule
Standard tactics & -- & -- & 14.0 & 14.0 & 14.0 & \checkmark \\
llmstep & 2.8B & Pythia & -- & -- & 27.9 & \checkmark \\
Kimina-Distill & 1.5B & Qw2.5M & 42.6 & 56.2 & -- & \checkmark \\
Kimina-Distill & 1.7B & Qw3 & -- & 73.0 & -- & \checkmark \\
Kimina-RL & 1.7B & Qw3 & -- & \underline{76.6} & -- & \checkmark \\
\midrule
\textbf{Ours (Qwen3.5-2B)} & \textbf{2B} & \textbf{Qw3.5} & \textbf{XX.X} & \textbf{XX.X} & \textbf{XX.X} & \checkmark \\
\textbf{Ours (Qwen3-1.7B)} & \textbf{1.7B} & \textbf{Qw3} & \textbf{XX.X} & \textbf{XX.X} & \textbf{XX.X} & \checkmark \\
\midrule
BFS-Prover-V2 & 7B & Qw2.5M & -- & -- & 82.4 & -- \\
DSP-V2 & 7B & DSP1.5 & -- & -- & 88.9 & -- \\
\bottomrule
\end{tabular}
\label{tab:main}
\end{table}

\subsection{Training Stage Ablation (Best Base Model)}

\begin{table}[h]
\centering
\caption{Contribution of each stage. MiniF2F-test pass@32.}
\begin{tabular}{lcc}
\toprule
\textbf{Stage} & \textbf{Data Added} & \textbf{pass@32} \\
\midrule
SFT (LeanDojo only) & 200K & XX.X\% \\
SFT (full) & 1.2M & XX.X\% \\
+ Expert Iter. R1 & +150K & XX.X\% \\
+ Expert Iter. R2 & +100K & XX.X\% \\
+ Expert Iter. R3 & +70K & XX.X\% \\
+ GRPO & -- & XX.X\% \\
+ DPO & 50K pos/neg & XX.X\% \\
\bottomrule
\end{tabular}
\label{tab:ablation}
\end{table}

\subsection{Latency}

\begin{table}[h]
\centering
\caption{On-device inference latency (single tactic, Q4\_K\_M).}
\begin{tabular}{lcccc}
\toprule
\textbf{Hardware} & \multicolumn{2}{c}{\textbf{Qwen3.5-2B}} & \multicolumn{2}{c}{\textbf{Qwen3-1.7B}} \\
 & Prompt & Gen & Prompt & Gen \\
\midrule
M3 Pro & $\sim$3ms & $\sim$0.4s & $\sim$2ms & $\sim$0.3s \\
M4 Pro & $\sim$2ms & $\sim$0.3s & $\sim$1.5ms & $\sim$0.2s \\
M5 & $\sim$2ms & $\sim$0.3s & $\sim$2ms & $\sim$0.25s \\
\bottomrule
\end{tabular}
\label{tab:latency}
\end{table}

At 8 candidates per goal ($\sim$0.3s each), a full beam takes $\sim$2.4s. With Pantograph at 3ms/tactic, inference dominates search time.

\subsection{Quantization}

\begin{table}[h]
\centering
\caption{Quantization effect on quality. MiniF2F-test pass@32.}
\begin{tabular}{lccccc}
\toprule
\textbf{Quant} & \multicolumn{2}{c}{\textbf{Qwen3.5-2B}} & \multicolumn{2}{c}{\textbf{Qwen3-1.7B}} \\
 & Size & pass@32 & Size & pass@32 \\
\midrule
F16 & 4.0 GB & XX.X\% & 3.4 GB & XX.X\% \\
Q8\_0 & 2.1 GB & XX.X\% & 1.7 GB & XX.X\% \\
Q4\_K\_M & 1.5 GB & XX.X\% & 1.0 GB & XX.X\% \\
\bottomrule
\end{tabular}
\label{tab:quant}
\end{table}

%----------------------------------------------------------------------
\section{Analysis}
\label{sec:analysis}

\subsection{Qwen3.5 GDN vs Standard Transformer}

The hybrid GDN architecture in Qwen3.5-2B presents a novel tradeoff for tactic prediction:

\begin{itemize}
\item \textbf{Potential advantage}: O(1) per-token inference in 75\% of layers (linear attention). For short outputs like tactics (10--50 tokens), the prompt processing cost dominates, and GDN's efficient state compression may help with long tactic states.
\item \textbf{Potential disadvantage}: GDN layers are more sensitive to quantization. The \texttt{ssm\_beta} and \texttt{ssm\_alpha} parameters control the linear recurrence dynamics---quantization errors here propagate multiplicatively through the sequence, unlike attention weights where errors are bounded by the softmax.
\item \textbf{Fine-tuning}: LoRA at bf16 (not QLoRA) is required to preserve GDN layer fidelity, increasing training memory requirements.
\end{itemize}

\subsection{What the Model Learns}

Compared to standard tactics, the trained model generates:
\begin{itemize}
\item \textbf{Lemma-specific tactics}: \texttt{exact Nat.add\_comm}, \texttt{rw [mul\_assoc]}
\item \textbf{Tactic combinations}: \texttt{simp [h1, h2]; linarith}
\item \textbf{Hypothesis-aware tactics}: \texttt{exact h0}, \texttt{rw [$\leftarrow$ h2]}
\end{itemize}

\subsection{Failure Modes}

\begin{itemize}
\item \textbf{Hallucinated lemma names}: Primary failure. The model generates plausible-sounding but non-existent Mathlib names. Caught by Pantograph in 3ms.
\item \textbf{Deep proofs ($>$10 steps)}: Still challenging. The length penalty helps but multi-step reasoning degrades with model size.
\item \textbf{Out-of-distribution}: Category theory and advanced topology see lower success rates.
\end{itemize}

\subsection{Expert Iteration Dynamics}

Diminishing returns in new proofs per round, but increasing proof quality (shorter proofs, more robust tactics). Round 1 discovers the most; rounds 2--3 refine.

%----------------------------------------------------------------------
\section{Discussion}
\label{sec:discussion}

\paragraph{On-device theorem proving.} A $\sim$1.5GB model producing sub-second tactics with no network dependency removes the main adoption barrier for proof assistants. Any Apple Silicon Mac can run it.

\paragraph{Architecture matters.} If Qwen3.5-2B outperforms Qwen3-1.7B despite having only 0.3B more parameters, it suggests that hybrid architectures are a free upgrade for tactic prediction. If not, it validates the conservative choice of standard transformers for domain-specific tasks.

\paragraph{Complementarity with agentic proving.} This model is the inner loop of OpenProof's larger architecture (LLM-driven planning, decomposition, repair). Better tactic search improves all modes.

\paragraph{Scaling path.} The pipeline scales to 7B (Qwen3.5-9B or Qwen3-8B) for cloud or high-RAM machines. A 7B variant should approach BFS-Prover-V2 levels (80\%+).

\paragraph{Continuous improvement.} Every proof completed by users generates new (state, tactic) pairs. Expert iteration runs continuously as the corpus grows.

%----------------------------------------------------------------------
\section{Training Cost}

\begin{table}[h]
\centering
\caption{Estimated compute cost (Thunder Compute, A100-80GB at \$0.78/hr).}
\begin{tabular}{lcc}
\toprule
\textbf{Stage} & \textbf{GPU Hours} & \textbf{Cost} \\
\midrule
SFT ($\times$2 bases) & 16 & \$12 \\
Expert Iteration (3 rounds $\times$ 2) & 400 & \$312 \\
DAPO ($\times$2) & 48 & \$37 \\
Eval \& misc & 30 & \$23 \\
\midrule
\textbf{Total} & \textbf{$\sim$494} & \textbf{$\sim$\$385} \\
\bottomrule
\end{tabular}
\label{tab:cost}
\end{table}

%----------------------------------------------------------------------
\section{Conclusion}

We compared two base models for on-device Lean 4 tactic prediction: Qwen3.5-2B (the newest hybrid GDN architecture) and Qwen3-1.7B (a proven standard transformer). Both were trained with the same three-stage pipeline (SFT, expert iteration, GRPO) for under \$400 total. The best variant achieves XX.X\% pass@32 on MiniF2F-test at $\sim$1.5GB quantized---a new Pareto optimum for on-device theorem proving. We release model weights, training code, and the full 1.2M-pair dataset.

%----------------------------------------------------------------------
\bibliographystyle{plainnat}

\begin{thebibliography}{99}

\bibitem[BFS-Prover-V1(2025)]{bfsprover1}
Xin et al.
\newblock BFS-Prover: Scalable Best-First Tree Search for LLM-based Automatic Theorem Proving.
\newblock \emph{arXiv:2502.03438}, February 2025.

\bibitem[BFS-Prover-V2(2025)]{bfsprover2}
ByteDance Seed.
\newblock BFS-Prover-V2: Multi-Turn Off-Policy RL and Multi-Agent Tree Search for Formal Theorem Proving.
\newblock \emph{arXiv:2509.06493}, September 2025.

\bibitem[DeepSeekMath(2024)]{deepseekmath}
Shao et al.
\newblock DeepSeekMath: Pushing the Limits of Mathematical Reasoning in Open Language Models.
\newblock \emph{arXiv:2402.03300}, February 2024.

\bibitem[DeepSeek-Prover-V1.5(2024)]{deepseekmathv15}
Xin et al.
\newblock DeepSeek-Prover-V1.5: Harnessing Proof Assistant Feedback for RL and MCTS.
\newblock \emph{arXiv:2408.08152}, August 2024.

\bibitem[DeepSeek-Prover-V2(2025)]{deepseekprover2}
DeepSeek.
\newblock DeepSeek-Prover-V2: Advancing Formal Mathematical Reasoning via RL for Subgoal Decomposition.
\newblock \emph{arXiv:2504.21801}, April 2025.

\bibitem[Delta-Prover(2025)]{deltaprover}
Delta-Prover: Agentic Theorem Proving with Reflective Decomposition.
\newblock \emph{arXiv:2507.15225}, July 2025.

\bibitem[DAPO(2025)]{dapo}
Yu et al.
\newblock DAPO: An Open-Source LLM Reinforcement Learning System.
\newblock \emph{arXiv:2503.14476}, 2025.

\bibitem[DPO(2023)]{dpo}
Rafailov et al.
\newblock Direct Preference Optimization.
\newblock \emph{NeurIPS 2023}.

\bibitem[Dr.\ GRPO(2025)]{drgrpo}
Liu et al.
\newblock Understanding R1-Zero-Like Training: A Critical Perspective.
\newblock \emph{arXiv:2503.20783}, 2025.

\bibitem[Expert Iteration(2017)]{expertiteration}
Anthony, Tian, and Barber.
\newblock Thinking Fast and Slow with Deep Learning and Tree Search.
\newblock \emph{NeurIPS 2017}.

\bibitem[Falcon-H1(2025)]{falconh1}
TII.
\newblock Falcon-H1: A Family of Hybrid-Head Language Models.
\newblock \emph{arXiv:2507.22448}, July 2025.

\bibitem[Goedel-Prover-V2(2025)]{goedel}
Lin et al.
\newblock Goedel-Prover: A Frontier Model for Open-Source Automated Theorem Proving.
\newblock \emph{arXiv:2502.07640 / ICLR 2026}.

\bibitem[HILBERT(2025)]{hilbert}
Apple ML Research.
\newblock HILBERT: A Prototype System for Solving Competition-Level Math Problems.
\newblock \emph{arXiv:2509.22819}, September 2025.

\bibitem[Kimina-Prover(2025)]{kimina}
AI-MO / Moonshot AI.
\newblock Kimina-Prover: Towards Scalable Mathematical Reasoning through Large-Scale RL.
\newblock \emph{arXiv:2504.11354}, April 2025.

\bibitem[LeanDojo(2023)]{leandojo}
Yang et al.
\newblock LeanDojo: Theorem Proving with Retrieval-Augmented Language Models.
\newblock \emph{NeurIPS 2023}.

\bibitem[Lean Workbook(2024)]{leanworkbook}
Ying et al.
\newblock Lean Workbook: A Large-Scale Lean Problem Set Formalized from Natural Language Math Problems.
\newblock \emph{arXiv:2406.03847}, June 2024.

\bibitem[llmstep(2023)]{llmstep}
Welleck and Saha.
\newblock llmstep: LLM proofstep suggestions in Lean.
\newblock \emph{arXiv:2310.18457}, October 2023.

\bibitem[MiniF2F(2022)]{minif2f}
Zheng et al.
\newblock MiniF2F: A Cross-System Benchmark for Formal Olympiad-Level Mathematics.
\newblock \emph{ICLR 2022}.

\bibitem[Pantograph(2025)]{pantograph}
Aniva et al.
\newblock Pantograph: A Machine-to-Machine Interaction Interface for Advanced Theorem Proving.
\newblock \emph{TACAS 2025}.

\bibitem[Prompt-Aug GRPO(2026)]{promptauggrpo}
Prompt Augmentation for GRPO.
\newblock \emph{arXiv:2602.03190}, February 2026.

\bibitem[Qwen2.5-Math(2024)]{qwen25math}
Yang et al.
\newblock Qwen2.5-Math Technical Report.
\newblock \emph{arXiv:2409.12122}, September 2024.

\bibitem[Qwen3(2025)]{qwen3}
Qwen Team.
\newblock Qwen3 Technical Report.
\newblock \emph{arXiv:2505.09388}, May 2025.

\bibitem[Qwen3.5(2026)]{qwen35}
Qwen Team.
\newblock Qwen3.5: Scaling Hybrid Attention to Small Models.
\newblock \emph{GitHub / HuggingFace release}, March 2026.

\bibitem[RLAD(2026)]{rlad}
Reinforcement Learning-Aware Knowledge Distillation.
\newblock \emph{arXiv:2602.22495}, February 2026.

\bibitem[rStar-Math(2025)]{rstarmath}
Guan et al.
\newblock rStar-Math: Small LLMs Can Master Math Reasoning with Self-Evolved Deep Thinking.
\newblock \emph{arXiv:2501.04519}, January 2025.

\bibitem[Seed-Prover 1.5(2025)]{seedprover15}
ByteDance Seed.
\newblock Seed-Prover 1.5.
\newblock \emph{arXiv:2512.17260}, December 2025.

\bibitem[TRL(2024)]{trl}
von Werra et al.
\newblock TRL: Transformer Reinforcement Learning.
\newblock \url{https://github.com/huggingface/trl}.

\bibitem[TTRL(2025)]{ttrl}
Test-Time Reinforcement Learning.
\newblock \emph{arXiv:2504.16084}, April 2025.

\bibitem[Unsloth(2024)]{unsloth}
Han and Han.
\newblock Unsloth: 2x Faster LLM Fine-tuning.
\newblock \url{https://github.com/unslothai/unsloth}.

\bibitem[V-STaR(2024)]{vstar}
Hosseini et al.
\newblock V-STaR: Training Verifiers for Self-Taught Reasoners.
\newblock \emph{arXiv:2402.06457}, February 2024.

\end{thebibliography}

\end{document}
